\section{第五章：Fourier 级数与 Fourier 变换}

\begin{frame}{第五章：为什么要做频谱分解？}
  \begin{itemize}
    \item 复杂信号往往由许多简单振动叠加而成。
    \item 频率分量可分别分析、过滤和传播。
    \item 微分方程在频域中常化为代数方程或常微分方程。
  \end{itemize}
  \begin{block}{核心问题}
    如何用一组正交的正弦、余弦或复指数函数表示给定函数？
  \end{block}
\end{frame}

\begin{frame}{函数空间中的内积与正交}
  在 $[a,b]$ 上定义内积
  \[
    (f,g)=\int_a^b f^*(x)g(x)\,\mathrm{d}x.
  \]
  若 $(f,g)=0$，称 $f$ 与 $g$ 正交；范数为 $\|f\|^2=(f,f)$。
  在 $[-\pi,\pi]$ 上，
  \[
    \int_{-\pi}^{\pi}\sin nx\sin mx\,\mathrm{d}x=\pi\delta_{mn},
    \qquad
    \int_{-\pi}^{\pi}\sin nx\cos mx\,\mathrm{d}x=0.
  \]
  \begin{block}{类比}
    Fourier 系数就像普通向量在正交坐标轴上的分量；正交性使每个频率分量能被单独投影出来。
  \end{block}
\end{frame}

\begin{frame}{Fourier 级数与系数}
  对周期为 $2\ell$ 的函数，写作
  \[
    f(x)=\frac{a_0}{2}+\sum_{n=1}^{\infty}
    \left[a_n\cos\left(\frac{n\pi x}{\ell}\right)+
    b_n\sin\left(\frac{n\pi x}{\ell}\right)\right].
  \]
  正交投影给出
  \[
    a_n=\frac{1}{\ell}\int_{-\ell}^{\ell}f(x)
    \cos\left(\frac{n\pi x}{\ell}\right)\mathrm{d}x,
    \quad
    b_n=\frac{1}{\ell}\int_{-\ell}^{\ell}f(x)
    \sin\left(\frac{n\pi x}{\ell}\right)\mathrm{d}x.
  \]
  \begin{itemize}
    \item 偶函数只含余弦项；奇函数只含正弦项。
    \item 在跳跃点，级数收敛到左右极限的平均值。
  \end{itemize}
\end{frame}

\begin{frame}{例：热方程为何自然产生 Fourier 级数？}
  对长度为 $\ell$ 的杆，零端点温度满足
  \[
    T_t=T_{xx},\qquad T(0,t)=T(\ell,t)=0.
  \]
  分离变量产生一组空间本征函数
  \[
    \sin\left(\frac{n\pi x}{\ell}\right),
  \]
  以及对应的时间衰减
  \[
    e^{-n^2\pi^2t/\ell^2}.
  \]
  所以
  \[
    T(x,t)=\sum_{n=1}^{\infty}A_n
    \sin\left(\frac{n\pi x}{\ell}\right)e^{-n^2\pi^2t/\ell^2}.
  \]
  \begin{alertblock}{物理读法}
    高频模态衰减得更快，因此扩散会优先抹平尖锐、细小的温度起伏。
  \end{alertblock}
\end{frame}

\begin{frame}{从 Fourier 级数到 Fourier 变换}
  当周期 $2\ell\to\infty$ 时，离散频率间隔趋于零，求和变为积分。采用一组常用约定：
  \[
    \hat f(k)=\int_{-\infty}^{\infty}f(x)e^{-\mathrm{i}kx}\,\mathrm{d}x,
    \qquad
    f(x)=\frac{1}{2\pi}\int_{-\infty}^{\infty}\hat f(k)e^{\mathrm{i}kx}\,\mathrm{d}k.
  \]
  \begin{itemize}
    \item $x$ 空间描述信号的位置或时间行为。
    \item $k$ 空间描述波数或频率成分。
    \item 微分变换为乘法：$\widehat{f'}(k)=\mathrm{i}k\hat f(k)$。
  \end{itemize}
\end{frame}

\begin{frame}{Delta 函数与 Laplace 变换}
  Dirac delta 满足筛选性质
  \[
    \int_{-\infty}^{\infty}\delta(x-x_0)f(x)\,\mathrm{d}x=f(x_0).
  \]
  它描述理想化的点源、瞬时冲量和 Green 函数。
  \begin{block}{Laplace 变换}
    \[
      F(s)=\mathcal{L}\{f(t)\}(s)=\int_0^{\infty}e^{-st}f(t)\,\mathrm{d}t.
    \]
    对初值问题，
    \[
      \mathcal{L}\{f'(t)\}=sF(s)-f(0).
    \]
  \end{block}
  Laplace 变换把时间微分方程转化为代数方程，同时自然保留初始条件。
\end{frame}
